1: % Reference entries, R (AN-11).  Law 1: os/iword/app/face_a.c
2: % retrieve_named (a missing file, a file that is not a WordPerfect 5.1
3: % document, a file larger than 64K; onto a clear screen the file becomes the
4: % document, into a document its text goes in at the cursor), os/iword/fs/
5: % wp5.c iw_wp5_read (only codes the corpus verified are accepted),
6: % spec/iword/wp5_format.h (HRt 0x0A, SRt 0x0D, DSRt 0x90), edit.c (wrap,
7: % the DSRt of Justification), codes.c iw_reveal_pane and face_a.c
8: % paint_reveal (the Reveal Codes screen: 10 document rows, the status line,
9: % the ruler, the 10-row codes area, the restore message; spec/iword/
10: % reveal_codes.h).  Shots: ../wordperfect51-decomp goldens MINT01..05,
11: % mint MINT06..32.
12: % Replace (Alt-F2) is P3 (bead initech-qvgd part 2, WL-0146): face_a.c
13: % IW_FK_replace, replace_key (w/Confirm? No (Yes): y confirms, n or Enter
14: % does not, other keys ignored, F1 cancels; Replace? Yes (No): y or Enter
15: % replaces, n skips, F1 stops), search_key in M_RSEARCH (the Srch: keys; F2
16: % / Esc go on to Replace with:, an empty string returns to the document) and
17: % M_RTEXT (no F6 / F8, no Up / Down), replace_next (every match from the
18: % cursor in the direction shown, then "* Not found *"), replace_match (a
19: % lower-case replacement takes the capitals of the text found: all capitals,
20: % or the first letter; an empty replacement deletes; Tab, Indent and Hard
21: % Page codes in the replacement get their positions fixed up).  F11 is
22: % Reveal Codes (spec/iword/fkeys.json; Reference p. 966 note 3).
23: % GATES: test-pubs-iword test-iword-face test-iword-wp5-roundtrip test-iword-reveal-codes test-flair-iword test-flair-iword-codes test-pubs-flair-screens test-iword-replace test-flair-iword-replace test-iword-fkeys
24: \Gates{test-pubs-iword test-iword-face test-iword-wp5-roundtrip test-iword-reveal-codes test-flair-iword test-flair-iword-codes test-pubs-flair-screens test-iword-replace test-flair-iword-replace test-iword-fkeys}
25: \begin{Feature}{Replace}
26: Replace finds text or codes, as Search does, and puts other text in their
27: place.  It can replace every occurrence at once, or stop at each one and ask
28: you.
29: 
30: To replace,
31: \begin{WPSteps}
32: \item Move the cursor to where the replacing is to begin: Replace works from
33: the cursor to the end of the document.  (To go to the top of the document,
34: press \K{Home, Home, Up}.)
35: \item Press \FK{Replace}{Alt-F2}.  The status line asks
36: \Scr{w/Confirm? No (Yes)}.
37: \item Type \Typed{y} to be asked at each occurrence
38: \Or
39: Type \Typed{n} or press \K{Enter} to replace every occurrence without
40: asking.
41: \item At \Scr{-> Srch:}, type the text to be replaced and press \K{F2}.
42: \item At \Scr{Replace with:}, type the new text and press \K{F2}.
43: \end{WPSteps}
44: \Documents{wpkey}{Alt-F2}
45: \Documents{wpkey}{Rplc:F2}
46: \Documents{wpkey}{Rplc:Esc}
47: InitechWord replaces the text wherever it finds it, from the cursor to the
48: end of the document, and leaves the cursor after the last replacement.  When
49: there is nothing more to find, the status line shows \Scr{* Not found *};
50: press any key to go on.
51: 
52: \BeginReplay{ref-replace}
53: For example, type
54: \Type{The memo of Friday replaces the memo of Monday.  File the MEMO.}
55: press \PressKey{Home, Home, Up} and \Press{Replace}{Alt-F2}, type
56: \TypeName{n}, type \TypeName{memo} and press \PressKey{F2}:
57: \WordScreen[rows=1]{ref-replace-1}{The replacement prompt}
58: Type \TypeName{notice} and press \PressKey{F2}:
59: \WordScreen[rows=2]{ref-replace-2}{Every memo replaced}
60: Each \textit{memo} is now \textit{notice}, and the capitals of \textit{MEMO}
61: are kept: \textit{NOTICE}.
62: \EndReplay
63: 
64: \Notes
65: \Topic{Canceling}
66: Press \FK{Cancel}{F1} at \Scr{w/Confirm? No (Yes)}, at the search string or
67: at \Scr{Replace with:} to leave Replace without replacing anything.  At
68: \Scr{Replace? Yes (No)}, Cancel stops the replacing; what was replaced
69: before it stays replaced.
70: \Documents{wpkey}{Rplc:F1}
71: 
72: \Topic{Capital Letters}
73: The search string finds capitals as described under \InRef{Search}.  When
74: the new text is all in lower case, InitechWord gives it the capitals of the
75: text it replaces: replacing \Typed{memo} with \Typed{notice} gives
76: \Scr{notice} for \Scr{memo}, \Scr{Notice} for \Scr{Memo} and \Scr{NOTICE}
77: for \Scr{MEMO}.  New text with any capital is put in exactly as you typed
78: it.
79: 
80: \Topic{Codes}
81: The search string can hold the codes Search can find (see \InRef{Search}).
82: Replacing \Code{[BOLD]} with nothing, for example, deletes every bold on
83: code; the \Code{[bold]} codes stay where they are, as Reveal Codes shows,
84: and change nothing.  To delete them too, replace \Code{[bold]} with
85: nothing: at the \Scr{-> Srch:} prompt press F6 twice, which gives
86: \Code{[BOLD][bold]}, then \K{Left} twice and \K{Delete}, which leaves
87: \Code{[bold]}.  The new text can hold \Code{[Tab]}
88: (\K{Tab}), \Code{[\rarr Indent]} (\K{F4}), \Code{[HRt]} (\K{Enter}) and
89: \Code{[HPg]} (\K{Ctrl-Enter}), but no appearance codes: F6 and F8 do nothing
90: at \Scr{Replace with:}.
91: \Documents{wpkey}{Rplc:Tab}
92: \Documents{wpkey}{Rplc:F4}
93: \Documents{wpkey}{Rplc:Enter}
94: \Documents{wpkey}{Rplc:Ctrl-Enter}
95: 
96: \Topic{Confirming Each Replacement}
97: If you answered \Typed{y} to \Scr{w/Confirm? No (Yes)}, InitechWord stops
98: at each occurrence, highlights it as a block and asks
99: \Scr{Replace? Yes (No)}.  Type \Typed{y} or press \K{Enter} to replace it, or
100: type \Typed{n} to leave it; InitechWord then goes on to the next.  Other
101: keys are ignored.
102: 
103: \BeginReplay{ref-replace-confirm}
104: For example, type
105: \Type{Send the report to Bill.  The report is due today.  Report early.}
106: press \PressKey{Home, Home, Up} and \Press{Replace}{Alt-F2}, type
107: \TypeName{y}, type \TypeName{report}, press \PressKey{F2}, type
108: \TypeName{memo} and press \PressKey{F2}:
109: \WordScreen[rows=1]{ref-replace-confirm-1}{Replace asks at each occurrence}
110: Type \TypeName{y} for the first, \TypeName{n} for the second and
111: \TypeName{y} for the third:
112: \WordScreen[rows=1]{ref-replace-confirm-2}{The first and third replaced}
113: \EndReplay
114: 
115: \Topic{Deleting Text Throughout}
116: To delete every occurrence of a word, replace it with nothing: at
117: \Scr{Replace with:} press F2 without typing anything.
118: 
119: \Topic{Direction}
120: Replace works forward, from the cursor to the end of the document.  To
121: replace backward, from the cursor to the beginning, press \K{Up} at the
122: \Scr{-> Srch:} prompt: it becomes \Scr{<- Srch:}.  \K{Down} turns it forward
123: again.
124: 
125: \Topic{Editing the Strings}
126: The search string and the new text are typed and corrected with the keys of
127: the \Scr{Srch:} prompt (see \InRef{Search}): \K{Left}, \K{Right},
128: \K{Backspace} and \K{Delete}, a whole code at a time.  \K{F2} or \K{Esc}
129: ends each string.  At \Scr{Replace with:} InitechWord does not look at
130: Shift, Alt or Ctrl on these keys:
131: \Documents{wpkey}{Rplc:Left}
132: \Documents{wpkey}{Rplc:Right}
133: \Documents{wpkey}{Rplc:Backspace}
134: \Documents{wpkey}{Rplc:Delete}
135: \par\addvspace{4pt}
136: \noindent{\small\begin{tabular}{@{}lll@{}}
137: \WPPromptAlias{Rplc:}{Shift-F1}{F1}, & \WPPromptAlias{Rplc:}{Alt-F1}{F1}, & \WPPromptAlias{Rplc:}{Ctrl-F1}{F1}\\
138: \WPPromptAlias{Rplc:}{Shift-F2}{F2}, & \WPPromptAlias{Rplc:}{Alt-F2}{F2}, & \WPPromptAlias{Rplc:}{Ctrl-F2}{F2}\\
139: \WPPromptAlias{Rplc:}{Shift-Esc}{Esc}, & \WPPromptAlias{Rplc:}{Alt-Esc}{Esc}, & \WPPromptAlias{Rplc:}{Ctrl-Esc}{Esc}\\
140: \WPPromptAlias{Rplc:}{Shift-Left}{Left}, & \WPPromptAlias{Rplc:}{Alt-Left}{Left}, & \WPPromptAlias{Rplc:}{Ctrl-Left}{Left}\\
141: \WPPromptAlias{Rplc:}{Shift-Right}{Right}, & \WPPromptAlias{Rplc:}{Alt-Right}{Right}, & \WPPromptAlias{Rplc:}{Ctrl-Right}{Right}\\
142: \WPPromptAlias{Rplc:}{Shift-Backspace}{Backspace}, & \WPPromptAlias{Rplc:}{Alt-Backspace}{Backspace}, & \WPPromptAlias{Rplc:}{Ctrl-Backspace}{Backspace}\\
143: \WPPromptAlias{Rplc:}{Shift-Delete}{Delete}, & \WPPromptAlias{Rplc:}{Alt-Delete}{Delete}, & \WPPromptAlias{Rplc:}{Ctrl-Delete}{Delete}\\
144: \WPPromptAlias{Rplc:}{Shift-Tab}{Tab}, & \WPPromptAlias{Rplc:}{Alt-Tab}{Tab}, & \WPPromptAlias{Rplc:}{Ctrl-Tab}{Tab}\\
145: \WPPromptAlias{Rplc:}{Shift-Enter}{Enter}, & \WPPromptAlias{Rplc:}{Alt-Enter}{Enter} & \\
146: \end{tabular}}
147: \par\addvspace{4pt}
148: 
149: \Topic{Repeating}
150: Replace starts with an empty search string.  Afterwards, the search string
151: of Replace is offered by the next Search (F2), so that you can look for any
152: occurrence you chose not to replace.
153: 
154: \Topic{Undoing a Replacement}
155: Replace does not keep what it replaced.  To go back, retrieve the document
156: as you last saved it, or replace the new text with the old.
157: 
158: \SeeAlso{Block; Codes; Search; Undelete}
159: \end{Feature}
160: 
161: \begin{Feature}{Retrieve}
162: Retrieve brings a document you have saved back onto the screen, so that you
163: can read it or change it.  You can also
164: retrieve a document into the one on the screen, to combine the two.
165: 
166: To retrieve a document,
167: \begin{WPSteps}
168: \item Clear the screen if it holds a document you do not want to combine
169: with the one you retrieve (see \InRef{Exit}).
170: \item Press \FK{Retrieve}{Shift-F10}.  The status line shows
171: \Scr{Document to be Retrieved:}.
172: \PullDown{Select \Mi{R}etrieve from the \Mi{F}ile menu.}
173: \item Type the name of the document and press \K{Enter}.
174: \end{WPSteps}
175: \Documents{wpkey}{Shift-F10}
176: On a clear screen the document you retrieve becomes the document you are
177: working on: its name is shown on the status line and the cursor is at its
178: top.
179: 
180: \Notes
181: \Topic{Into Another Document}
182: If the screen already holds text, the retrieved document's text is put into
183: it at the cursor, and the document on the screen keeps its own name.  The
184: two documents are now one; the file you retrieved is not changed.
185: 
186: \BeginReplay{ref-retrieve}
187: For example, with a new document on the screen, type
188: \TypeRun{Attached is the letter: } (with the space), press
189: \Press{Retrieve}{Shift-F10}, type \TypeName{letter} -- the letter saved in
190: \textit{A Brief Lesson} -- and press \PressKey{Enter}:
191: \WordScreen[rows=6]{ref-retrieve}{A document retrieved into another, at the cursor}
192: \EndReplay
193: The text of the letter follows the text you typed, and the status line shows
194: no name: the document on the screen has not been saved.
195: 
196: \Topic{Messages}
197: \Scr{ERROR: File Not Found} means there is no document by that name on the
198: data volume: check the name and type it again.
199: \Scr{ERROR: Incompatible File Format} means the file is not a WordPerfect
200: 5.1 document that InitechWord can read.
201: \Scr{ERROR: Document too large} means the file is larger than InitechWord
202: can hold (see Appendix H).  Press any key to return to the prompt.
203: 
204: \Topic{WordPerfect 5.1 Documents}
205: InitechWord's documents are WordPerfect 5.1 documents, and InitechWord
206: retrieves documents written by WordPerfect 5.1 itself, keeping every code
207: it holds (see \InRef{Codes}).  A WordPerfect document that holds a code
208: InitechWord does not know is refused with \Scr{ERROR: Incompatible File
209: Format}, rather than changed.
210: 
211: \SeeAlso{Exit; Filenames; Save}
212: \end{Feature}
213: 
214: \begin{Feature}{Return, Deletable Soft}
215: A deletable soft return ends a line without ending the paragraph.
216: InitechWord puts one in when you set the justification with the cursor in
217: the middle of a line, so that the text after the cursor starts a new line
218: in the new justification (see \InRef{Justification}).
219: 
220: Reveal Codes shows it as \Code{[DSRt]}.  To join the two lines again, delete
221: the code: InitechWord deletes it without asking.
222: 
223: \SeeAlso{Justification; Return, Hard; Return, Soft}
224: \end{Feature}
225: 
226: \begin{Feature}{Return, Hard}
227: A hard return ends a line where you choose, and ends the paragraph.  You put
228: one in by pressing \K{Enter}.  Use it at the end of a paragraph, after each
229: line of an address or a list, and to make blank lines.
230: 
231: To insert a hard return,
232: \begin{WPSteps}
233: \item Press \K{Enter}.
234: \end{WPSteps}
235: \Documents{wpkey}{Enter}
236: 
237: \Notes
238: \Topic{Code}
239: A hard return is the code \Code{[HRt]}.  Delete it to join a line to the
240: line after it (see \InRef{Delete}).
241: 
242: \Topic{Shift-Enter and Alt-Enter}
243: Shift-Enter and Alt-Enter put in a hard return as Enter does.  Ctrl-Enter
244: puts in a hard page (see \InRef{Page Break, Hard}).
245: 
246: \SeeAlso{Page Break, Hard; Return, Soft; Word Wrap}
247: \end{Feature}
248: 
249: \begin{Feature}{Return, Soft}
250: A soft return is the end of a line that InitechWord makes for you when a
251: line is full (see \InRef{Word Wrap}).  It takes the place of the space
252: between the last word that fits on the line and the word that goes to the
253: next line.  Reveal Codes shows it as \Code{[SRt]}.
254: 
255: You never type a soft return and you need not delete one: whenever you
256: change a paragraph, InitechWord wraps it again and moves its soft returns to
257: where the lines now end.
258: 
259: \SeeAlso{Return, Deletable Soft; Return, Hard; Word Wrap}
260: \end{Feature}
261: 
262: \begin{Feature}{Reveal Codes}
263: Reveal Codes shows the hidden codes of the document.  The screen divides:
264: the upper part shows the document as usual, and the lower part shows the
265: same lines with every code in square brackets.  Use it to see why text
266: looks as it does, and to delete a code.
267: 
268: To reveal codes,
269: \begin{WPSteps}
270: \item Press \FK{Reveal Codes}{Alt-F3}.
271: \item Move the cursor as usual.  It moves in both parts of the screen; in
272: the codes area it is the highlighted character or code.
273: \item Press \FK{Reveal Codes}{Alt-F3} again to return to the full editing
274: screen.
275: \end{WPSteps}
276: \Documents{wpkey}{Alt-F3}
277: \Documents{wpkey}{F11}
278: \K{F11} does what Alt-F3 does, with Block on or off.
279: \FlairShot{reveal}{640}{440}{Reveal Codes: the bold codes around a word}
280: 
281: The Reveal Codes screen has five parts, from the top:
282: \begin{description}
283: \item[The document] ten lines of the document, as on the editing screen.
284: \item[The status line] as on the editing screen.
285: \item[The ruler] the margins and tab stops in effect at the cursor.
286: \item[The codes area] ten rows: the text of the lines around the cursor,
287: starting three lines above it, with the codes shown.  A line longer than the
288: screen runs on into the next row.
289: \item[The message] \Scr{Press Reveal Codes to restore screen}.
290: \end{description}
291: 
292: \Notes
293: \Topic{Deleting Codes}
294: In Reveal Codes, Backspace and Delete delete codes without asking, as they
295: delete characters.  Deleting either code of an appearance or size deletes
296: its partner too.  See \InRef{Codes}.
297: 
298: \Topic{Editing in Reveal Codes}
299: Everything you can do on the editing screen you can do in Reveal Codes:
300: type, block, move, format.  Both parts of the screen change as you work.
301: 
302: \Topic{The Cursor in the Codes Area}
303: The highlight covers the character or the whole code the cursor is on.  At
304: the very end of the document it is one empty cell after the last code.
305: 
306: \Topic{The End of the Document}
307: \BeginReplay{ref-reveal-end}At the end of the document the highlight is one empty cell after the
308: last character or code.  For example, type \TypeRun{Total:}, press
309: \PressKey{Tab}, type \TypeRun{570} and press \Press{Reveal Codes}{Alt-F3}:
310: \WordScreen[rows=13]{ref-reveal-end}{The cursor at the end of the document, in Reveal Codes}
311: \EndReplay
312: 
313: \Topic{The Ruler}
314: The left margin is shown by \Scr{\{} and the right margin by \Scr{\}}; a
315: margin that does not fall on a tab stop is shown by \Scr{[} or \Scr{]}
316: instead.  The small triangles are the tab stops, every half inch from the
317: left margin.
318: 
319: \SeeAlso{Codes; Margins, Left and Right; Tab}
320: \end{Feature}
